\chapter{Credit Relative Value}\label{ch:s2:credit-relative-value}

A company's bond and its credit default swap insure the same default. Buying the bond and buying protection should earn nothing but the risk-free
rate, so the basis, the CDS spread minus the bond's spread, should be zero. Bai and Collin-Dufresne found it deviating most for bonds with the
highest frictions: illiquid, expensive to fund, exposed to counterparties, poor as collateral. In 2008 the basis went deeply negative when the trade
that should close it needed funding nobody had. On this chapter's synthetic market the basis sits at $-15$ basis points until a funding crisis takes
it to $-151$. A negative-basis book levered ten times loses 60\% of its capital on the way down and, if it survives, recovers 76\% after the low.
Entered at the low, the same trade earns 70\% in a year. The build is \texttt{firm.creditrv}.

\section{The bond-CDS basis}

A bond's spread (Book~2, chapter~21, the z-spread over the risk-free curve) and a credit default swap on the same issuer (Book~2, chapter~23) both
price the issuer's default. A long bond hedged with CDS protection has no default risk: on a credit event the protection pays par less the bond's
recovery. In frictionless markets the bond spread equals the CDS spread. In practice the bond must be financed, the CDS posts collateral, and the
two sides are traded by different investors.

\begin{definition}[Negative basis trade]\label{def:s2:credit-relative-value:negbasis}
A \emph{negative basis trade}\index{negative basis trade} buys a bond whose spread exceeds its issuer's CDS spread, finances it in repo, and buys CDS
protection on the same issuer; it earns the gap between the two spreads less its funding, holds no default risk, and loses on its marks when the gap
widens further.
\end{definition}

\texttt{firm.creditrv} builds forty synthetic issuers with CDS spreads near 120 basis points, averaging 136 over the sample. Each bond's spread is its
CDS spread plus a funding premium common to all bonds (15 basis points in normal times) plus noise. The basis is therefore minus the funding premium.
In year 5 a planted crisis lifts the premium to 150 basis points over a quarter and brings it back over a year
(\cref{fig:s2:credit-relative-value:basis}). The average basis falls from $-14.9$ to $-150.7$ basis points, at year 5.25.

\section{Index against constituents}

\begin{definition}[Credit index arbitrage]\label{def:s2:credit-relative-value:index}
\emph{Credit index arbitrage}\index{credit index arbitrage} trades a CDS index against the basket of its constituents' single-name CDS when the index spread differs from the
average of the members' spreads (the index skew): selling protection on the cheaper side and buying it on the dearer.
\end{definition}

The index trades more than its members and moves first. The gap between them, the skew, reflects that liquidity, and it closes as the names catch up.
The synthetic skew is a mean-reverting noise of 4 basis points with a 20-day half-life. Trading it once it passes one standard deviation, and holding
until it crosses zero, earned 55 basis points of notional a year with a Sharpe ratio of 0.97. The book was in a position 60\% of the time and paid 1
basis point of spread (4.5 of notional) to trade each unit, since the single-name leg is costly.

\section{Credit curve trades}

\begin{definition}[Credit curve trade]\label{def:s2:credit-relative-value:curve}
A \emph{credit curve trade}\index{credit curve trade} buys protection at one maturity and sells it at another on the same issuer or index, in amounts that offset their spread
durations, betting on the slope of the credit curve rather than its level.
\end{definition}

A flattener sells protection at the short end and buys it at the long end, gaining when the long spread falls relative to the short one. It carries
default risk between the maturities: a default before the short contract matures pays both legs, but one after it pays only the long leg. The synthetic
ten-year minus five-year slope averages 40 basis points and reverts with a 60-day half-life. The banded trade earned 49 basis points a year, a Sharpe
ratio of 0.75, in a position 62\% of the time.

\section{Funding and the negative basis}

The negative-basis book holds every bond against protection on its issuer, finances the bonds with a 10\% haircut, and pays its own funding spread on
the financed part (\cref{lst:s2:credit-relative-value:book}). Capital is the haircut, so the book is levered ten times.

\begin{center}
\small
\begin{tabular}{@{}lcccc@{}}
\toprule
on capital, ten times & before (\%/yr) & to the low (\%) & after the low (\%) & ten years (\%)\\
\midrule
own funding 0 bp & 1.6 & $-59.9$ & 76.0 & 24.4\\
own funding 30 bp & $-1.1$ & $-60.6$ & 63.2 & $-2.6$\\
own funding 60 bp & $-3.8$ & $-61.2$ & 50.4 & $-29.6$\\
\bottomrule
\end{tabular}
\end{center}

Before the crisis the trade barely pays. A basis of $-15$ basis points covers little funding, and a trader who funds at 30 basis points loses money
holding it. When the crisis widens the basis by 135 basis points, a spread duration of 4.5 turns that into a 6.1\% loss on the notional: 61\% of the
capital at ten times (\cref{fig:s2:credit-relative-value:basis}). Nothing defaulted and the hedge worked. The loss is marks, and a book that must meet
margin calls sells at the low and never sees the recovery. Mitchell and Pulvino described how that happened in 2008. The imminent failure of prime
brokers cut hedge funds' leverage suddenly, and arbitrageurs who should have bought cheap bonds had to sell them.

The other side of the crisis is the best trade in the book. Entered at the basis's low and held a year, the negative-basis book earned 69.8\% of its
capital with no funding cost and 67.1\% at 30 basis points. Carry provided 8.4 and 5.7 points of that, and the basis's return to normal the other
61.4. The capital that can buy at the low is the capital that was not levered to the limit before it.

\begin{omfigure}
\begin{tikzpicture}
\begin{axis}[width=11cm,height=5.4cm,xlabel={\small year},ylabel={\small bp or \% of capital},tick label style={font=\footnotesize},
  xmin=0,xmax=10,ymin=-160,ymax=40,grid=major,grid style={gray!25},table/col sep=comma,
  legend style={font=\scriptsize,at={(0.5,-0.3)},anchor=north},legend columns=3]
\addplot[black,thick] table[x=year,y=basis]{figdata/strategies-2/18-credit-relative-value/basis.csv};
\addplot[omThm,thick] table[x=year,y=book0]{figdata/strategies-2/18-credit-relative-value/basis.csv};
\addplot[omDef,thick,dashed] table[x=year,y=book30]{figdata/strategies-2/18-credit-relative-value/basis.csv};
\legend{average basis (bp),book at 0 bp funding (\%),book at 30 bp funding (\%)}
\end{axis}
\end{tikzpicture}

\omcaption{The synthetic bond-CDS basis over ten years, with a funding crisis in year 5, and the cumulative return on capital of a negative-basis
book held from the start, ten times levered, for two funding costs. Data: \texttt{s2\_creditrv.market}.}\label{fig:s2:credit-relative-value:basis}
\end{omfigure}

\section{Strategy files}

\begin{strategyfile}{Negative basis trade}\label{strat:s2:credit-relative-value:negbasis}
\sfield{Who pays you, and why} Holders forced to sell bonds when funding is scarce; the basis pays for the balance sheet that holds them.
\sfield{Instruments and venues} Corporate bonds in repo; single-name CDS.
\sfield{Signal} The basis against one's own funding and haircut.
\sfield{Sizing and execution} Leverage set for a crisis widening, not the normal basis; term funding.
\sfield{Costs} Funding, haircuts, CDS spreads and collateral.
\sfield{How it dies} A funding crisis that widens the basis while leverage is withdrawn.
\sfield{Horizon, capacity, infrastructure} Months to years; repo lines and CDS documentation.
\sfield{Backtest honestly} Funding and haircuts as they changed; forced sales at the low.
\sfield{Sources} Bai and Collin-Dufresne (2019); Mitchell and Pulvino (2012); this chapter: $-60\%$ of capital at ten times in the crisis, $+70\%$ in the
year after its low.
\end{strategyfile}

\begin{strategyfile}{Index versus single-name arbitrage}\label{strat:s2:credit-relative-value:index}
\sfield{Who pays you, and why} Index traders who move the index before the single names.
\sfield{Instruments and venues} CDS indices; the members' single-name CDS.
\sfield{Signal} The index skew against its history.
\sfield{Sizing and execution} Banded entry, exit on convergence.
\sfield{Costs} Single-name spreads, the largest cost.
\sfield{How it dies} Illiquid names; members' credit events between rolls.
\sfield{Horizon, capacity, infrastructure} Weeks; a basket execution system.
\sfield{Backtest honestly} Single-name costs for the whole basket.
\sfield{Sources} This chapter: 55 basis points a year, Sharpe ratio 0.97 (planted skew).
\end{strategyfile}

\begin{strategyfile}{Credit curve flattener}\label{strat:s2:credit-relative-value:curve}
\sfield{Who pays you, and why} Demand for protection at one maturity (short-dated hedging in stress, long-dated from structured products).
\sfield{Instruments and venues} CDS at two maturities, on names or indices.
\sfield{Signal} The slope against its history and the issuer's credit quality.
\sfield{Sizing and execution} Spread-duration matched.
\sfield{Costs} Two legs.
\sfield{How it dies} Default between the maturities; a curve that inverts in distress.
\sfield{Horizon, capacity, infrastructure} Months.
\sfield{Backtest honestly} Jump-to-default on the mismatched leg.
\sfield{Sources} This chapter: 49 basis points a year, Sharpe ratio 0.75 (planted slope).
\end{strategyfile}

\begin{strategyfile}{Senior versus subordinated}\label{strat:s2:credit-relative-value:seniority}
\sfield{Who pays you, and why} Mispricing between claims of different seniority on one issuer.
\sfield{Instruments and venues} Senior and subordinated bonds or CDS.
\sfield{Signal} The ratio of their spreads against implied recovery assumptions.
\sfield{Sizing and execution} Hedged to the issuer's default risk.
\sfield{Costs} Illiquid subordinated paper.
\sfield{How it dies} Bail-in or restructuring rules that change recoveries.
\sfield{Horizon, capacity, infrastructure} Months.
\sfield{Backtest honestly} Recoveries from actual resolutions.
\sfield{Sources} No performance figure verified.
\end{strategyfile}

\begin{strategyfile}{Cross-currency credit}\label{strat:s2:credit-relative-value:currency}
\sfield{Who pays you, and why} Issuers' bonds priced differently in two currencies after swapping.
\sfield{Instruments and venues} Bonds of one issuer in two currencies; cross-currency swaps.
\sfield{Signal} The swapped spread difference.
\sfield{Sizing and execution} Currency and rate risk swapped away.
\sfield{Costs} The cross-currency basis (chapter~14).
\sfield{How it dies} The basis itself moves.
\sfield{Horizon, capacity, infrastructure} Months; swap lines.
\sfield{Backtest honestly} Swap costs at the time.
\sfield{Sources} No performance figure verified.
\end{strategyfile}

\section{Tutorial: same default, two prices again}\label{tut:s2:credit-relative-value:tut}

\begin{tutorial}
\textbf{Goal.} Simulate bonds and CDS with a funding premium, run the negative-basis book through a funding crisis, and trade the index skew and the
curve. \textbf{End state:} the table and the figure.
\begin{tutsteps}
\item \textbf{The books}.
\omcode{code/firm/creditrv/firm_creditrv.py}{84}{111}{The negative-basis book and the banded relative-value trades.}\label{lst:s2:credit-relative-value:book}
\item \textbf{The crisis}.
\omcode{code/strategies-2/18-credit-relative-value/python/s2_creditrv.py}{46}{66}{The book through the crisis, and the trade entered at the low.}\label{lst:s2:credit-relative-value:crisis}
\item \textbf{Run} \texttt{basis\_stats()}, \texttt{basis\_book()}, \texttt{at\_the\_trough()}, \texttt{rv\_trades()} and \texttt{fig\_creditrv.py}.
\end{tutsteps}
\textbf{What to change next.} Add a margin rule that cuts the book when capital falls below its haircut and measure the loss locked in; make the
funding premium differ by bond liquidity; add a default.
\end{tutorial}

\section{Build: credit relative value}\label{bld:s2:credit-relative-value:creditrv}

\begin{build}
\bfield{Purpose} A bond-CDS basis driven by funding, the negative-basis book, index-constituent and curve trades.
\bfield{Interface} \texttt{CreditConfig(\dots)}, \texttt{simulate\_credit(cfg)}, \texttt{negative\_basis(sim, cfg, own\_funding)},
\texttt{index\_arbitrage(sim, cfg)}, \texttt{curve\_trade(sim, cfg)}.
\bfield{Rules} The basis is minus the funding premium plus noise; marks at the spread duration; banded trades with costs per unit traded.
\bfield{Acceptance tests} \texttt{code/firm/creditrv/tests/}: the basis equals minus the funding premium without noise; carry and marks by hand; the
banded position rule.
\bfield{Stretch} Defaults; margin calls; liquidity-dependent funding.
\end{build}

\begin{omsources}
\item J.~Bai and P.~Collin-Dufresne, ``The CDS-bond basis'', \emph{Financial Management} 48(2), 2019.
\item M.~Mitchell and T.~Pulvino, ``Arbitrage crashes and the speed of capital'', \emph{Journal of Financial Economics} 104(3), 2012.
\end{omsources}

\section{Exercises}

\begin{exercise}[$\star$]\label{exo:s2:credit-relative-value:1}
The basis is $-15$ basis points, a trader funds at 30 over the risk-free rate, and the haircut is 10\%. What is the carry on capital, a year?
\end{exercise}

\begin{exercise}[$\star$]\label{exo:s2:credit-relative-value:2}
The basis widens from $-15$ to $-150$ basis points with a spread duration of 4.5. What does a book levered ten times lose on capital?
\end{exercise}

\begin{exercise}[$\star$]\label{exo:s2:credit-relative-value:3}
What is the highest leverage a 10\% haircut allows?
\end{exercise}

\begin{exercise}[$\star\star$]\label{exo:s2:credit-relative-value:4}
Why is the negative-basis trade free of default risk but not of loss?
\end{exercise}

\begin{exercise}[$\star\star$]\label{exo:s2:credit-relative-value:5}
Why did Bai and Collin-Dufresne find the basis more negative when the bond's lending fee is high?
\end{exercise}

\begin{exercise}[$\star\star$]\label{exo:s2:credit-relative-value:6}
What default risk does a credit curve flattener keep?
\end{exercise}

\begin{exercise}[$\star\star\star$]\label{exo:s2:credit-relative-value:7}
\emph{Coding.} Add a rule that sells the negative-basis book whenever its cumulative loss reaches 50\% of capital. What does it lock in, and what
does it miss?
\end{exercise}

\begin{exercise}[$\star\star\star$]\label{exo:s2:credit-relative-value:8}
\emph{Find the flaw.} ``The negative basis is an arbitrage: we hold to maturity and cannot lose.''
\end{exercise}

\section{Problem: Same Default, Two Prices Again}

\begin{problem}[{Weekend problem --- credit relative value}]\label{pb:s2:credit-relative-value:1}
The chapter's synthetic issuers and the public record.

\medskip\noindent\textbf{Part I --- The basis.}
\begin{enumerate}
\item Define the bond-CDS basis and the \omterm{def:s2:credit-relative-value:negbasis}{negative basis trade}.
\item Why should the basis be zero, and why is it not?
\item What did Bai and Collin-Dufresne find?
\item Describe the synthetic funding premium and crisis.
\end{enumerate}

\medskip\noindent\textbf{Part II --- Other trades.}
\begin{enumerate}[resume]
\item Define \omterm{def:s2:credit-relative-value:index}{credit index arbitrage} and give its synthetic result.
\item Define a \omterm{def:s2:credit-relative-value:curve}{credit curve trade} and give its synthetic result.
\item What risks do these two trades carry?
\item Why is the single-name leg the costly one?
\end{enumerate}

\medskip\noindent\textbf{Part III --- The crisis.}
\begin{enumerate}[resume]
\item Give the book's returns before the crisis for each funding cost.
\item Give the loss to the low and the recovery.
\item What did Mitchell and Pulvino describe?
\item What did the trade earn when entered at the low?
\end{enumerate}

\medskip\noindent\textbf{Part IV --- The verdict.}
\begin{enumerate}[resume]
\item State the \emph{named result}: the negative-basis trade's carry and its mark-to-market loss when funding costs jump.
\item Why does the trade barely pay in normal times?
\item Who can buy at the low?
\item How would you size a negative-basis book?
\item How would you backtest it honestly?
\item Which strategy file depends on the cross-currency basis?
\item How does this chapter relate to chapters~8 and 11?
\item In one sentence: what does the negative basis price?
\end{enumerate}
\end{problem}

\section{Interview questions}

\begin{interviewq}[$\star$ \iqroles{trader}]\label{iq:s2:credit-relative-value:1}
What is the CDS-bond basis, and what does a negative basis mean?
\end{interviewq}

\begin{interviewq}[$\star\star$ \iqroles{trader}]\label{iq:s2:credit-relative-value:2}
Walk through a \omterm{def:s2:credit-relative-value:negbasis}{negative basis trade} and its cash flows, including a credit event.
\end{interviewq}

\begin{interviewq}[$\star\star$ \iqroles{risk}]\label{iq:s2:credit-relative-value:3}
A negative-basis book is levered ten times. What would you stress?
\end{interviewq}

\begin{interviewq}[$\star\star$ \iqroles{researcher}]\label{iq:s2:credit-relative-value:4}
How would you compute a CDS index's intrinsic value, and why does it differ from the index?
\end{interviewq}

\begin{interviewq}[$\star\star$ \iqroles{developer}]\label{iq:s2:credit-relative-value:5}
What reference data does a CDS book need, and what goes wrong at index rolls?
\end{interviewq}

\begin{interviewq}[$\star\star\star$ \iqroles{researcher}]\label{iq:s2:credit-relative-value:6}
Show that a bond hedged with CDS protection to maturity has, ignoring funding and the delivery option, the cash flows of a risk-free bond plus the
spread difference, and name what breaks the equivalence in practice.
\end{interviewq}
